% =============================================================================
% CHAPITRE 3 -- DEVELOPPEMENT PROGRESSIF DES TRAITEMENTS PYTHON
% =============================================================================

\chapter{Développement des scripts Python}
\label{chap:developpement_validation}

% =============================================================================
% 3.1 ORGANISATION GENERALE DES DEVELOPPEMENTS
% =============================================================================

\section{Organisation générale des développements}
\label{sec:organisation_developpements}

% -----------------------------------------------------------------------------
% 3.1.1 Une progression par formats de données
% -----------------------------------------------------------------------------

\subsection{Une progression par formats de données}
\label{subsec:progression_sources}

Le développement des nouveaux traitements Python ne s'est pas fait en une seule
fois. Il a été réalisé progressivement, format par format en fonction des
priorités du produit et des difficultés rencontrées.

Cette organisation permettait de travailler par étapes. Pour chaque format, il
fallait comprendre les données disponibles, identifier les métadonnées utiles,
développer le traitement adapté, puis vérifier les résultats obtenus.

Le travail a d'abord porté sur des formats qui étaient les plus utilisés par les clients de l'entreprise, comme les
fichiers rasters et vectoriels. Il s'est ensuite poursuivi avec des formats ou
technologies plus spécifiques : GeoPackage, données CAO, Oracle, GeoParquet,
SQL Server et, plus récemment, Nuage des points.

Cette progression correspondait aussi à la manière dont la migration devait être
menée dans Scan. Chaque ajout devait rester compatible avec le fonctionnement
existant du produit. Il ne s'agissait donc pas seulement de lire une donnée, mais
de produire une sortie exploitable par les autres composants.

% Image à inclure :
% Une frise simple montrant l'ordre de prise en charge des formats :
% rasters/vecteurs, GeoPackage, CAO, Oracle, GeoParquet, SQL Server,
% Nuage des points.

% -----------------------------------------------------------------------------
% 3.1.2 Périmètre fonctionnel selon les formats
% -----------------------------------------------------------------------------

\subsection{Périmètre fonctionnel selon les formats}
\label{subsec:perimetre_sources}

Le périmètre n'était pas identique pour tous les formats. Certains ont fait
l'objet d'un travail de documentation des métadonnées. D'autres ont nécessité
des traitements complémentaires, comme le listing ou la signature.

Pour les fichiers rasters, vectoriels et CAO, le travail réalisé concernait
uniquement la documentation. L'objectif était d'extraire les informations
techniques nécessaires à la création ou à la mise à jour des fiches de
métadonnées dans Scan. Le listing et la signature n'ont pas été développés pour
ces formats, car ils pouvaient être réalisés directement par le client de Scan.

Pour les formats suivants, le périmètre dépendait des besoins du produit, de la
nature des données et des possibilités offertes par les bibliothèques utilisées.
Cette approche progressive permettait d'avancer sans figer trop tôt une solution
unique qui n'aurait pas forcément convenu à toutes les technologies.

\begin{table}[H]
\centering
\begin{tabularx}{\textwidth}{>{\bfseries}p{3.2cm}X}
\toprule
Format ou technologie & Travail réalisé ou en cours \\
\midrule
Rasters, vecteurs et CAO & Documentation. \\
GeoPackage & Inventaire, documentation et signature. \\
Oracle & Inventaire, documentation et signature. \\
GeoParquet & Documentation. \\
SQL Server & Inventaire, documentation et signature. \\
Nuage des points & Travaux en cours sur l'inventaire, la documentation et la signature. \\
\bottomrule
\end{tabularx}
\caption{Progression des développements selon les formats étudiés}
\label{tab:progression_sources_developpement}
\end{table}

% =============================================================================
% 3.2 PREMIERS DEVELOPPEMENTS SUR LES FICHIERS GEOGRAPHIQUES
% =============================================================================

\section{Développements sur les fichiers géographiques}
\label{sec:premiers_developpements_fichiers}

% -----------------------------------------------------------------------------
% 3.2.1 Données rasters et vectorielles
% -----------------------------------------------------------------------------

\subsection{Données rasters et vectorielles}
\label{subsec:donnees_rasters_vecteurs}

Les premiers développements ont porté sur les données rasters et vectorielles.
Ces formats constituaient un bon point d'entrée, car ils permettaient de
travailler sur des formats fréquents dans les usages géomatiques.

Pour les données vectorielles, deux formats ont été pris en compte dans un
premier temps : le Shapefile et le GeoJSON. Ces formats sont très utilisés pour
échanger des couches géographiques. Le traitement devait donc être capable de
lire la couche, d'identifier sa géométrie, de récupérer les champs attributaires
et de produire les métadonnées attendues par Scan.

Pour les données rasters, plusieurs formats ont été pris en charge, notamment
\texttt{ECW}, GeoTIFF, \texttt{JP2}, \texttt{JPG}, \texttt{PNG} et d'autres
formats courants. Les informations à extraire n'étaient pas exactement les mêmes
que pour les données vectorielles. Il fallait par exemple tenir compte de la
taille de l'image, du nombre de bandes, du système de coordonnées et de
l'emprise.

Pour cette première partie, seule la documentation (lookup) a été développée.
Le but était de produire une description fiable de la ressource analysée, sans
développer le listing ni la signature, qui pouvaient déjà être réalisés par le
client de Scan pour ces formats.

GDAL et OGR ont été utilisés pour accéder à ces informations. Ces bibliothèques
offrent une interface commune pour lire de nombreux formats. Elles permettent
donc de limiter le nombre de traitements spécifiques à écrire pour chaque type
de fichier.

Ce premier travail a permis de mettre en place une base de développement. Il a
aussi permis d'identifier les premières différences de comportement entre les
formats, notamment sur la manière dont les systèmes de coordonnées, les champs
ou l'enveloppe convexe sont exposés.

% Image à inclure :
% Une capture d'un résultat de documentation sur un format raster ou vectoriel.

% -----------------------------------------------------------------------------
% 3.2.2 GeoPackage
% -----------------------------------------------------------------------------

\subsection{GeoPackage}
\label{subsec:geopackage}

Le travail s'est ensuite poursuivi avec le format GeoPackage. Il s'agit d'un
format de fichier permettant de stocker plusieurs couches géographiques dans un
même fichier. Il est souvent utilisé pour échanger des données tout en conservant
une structure proche d'une petite base de données.

La prise en charge de GeoPackage nécessitait donc de parcourir les couches
présentes dans le fichier, puis de produire les informations attendues pour
chacune d'elles. Le traitement devait rester cohérent avec la logique de Scan :
une ressource identifiée devait pouvoir être documentée de manière claire et
stable.

Le développement a porté sur les trois traitements attendus : l'inventaire, la
documentation et la signature. L'inventaire permettait de lister les couches
présentes dans le fichier. La documentation produisait les métadonnées de chaque
couche. La signature permettait ensuite de détecter les changements entre deux
analyses.

Ce format demandait une attention particulière, car un seul fichier peut contenir
plusieurs couches. Il fallait donc bien distinguer la ressource globale du
fichier et les couches réellement exploitées par Scan. Les résultats ont été
vérifiés sur plusieurs jeux de données afin de contrôler la cohérence des
sorties.

GeoPackage a aussi permis de confirmer l'importance d'une architecture souple.
Même si le format est stocké dans un fichier, son fonctionnement se rapproche
sur certains points d'une base de données. Le code devait donc éviter de traiter
tous les fichiers de manière trop uniforme.

% -----------------------------------------------------------------------------
% 3.2.3 GeoParquet
% -----------------------------------------------------------------------------

\subsection{GeoParquet}
\label{subsec:geoparquet}

GeoParquet a ensuite été étudié comme format de stockage géographique plus
récent. Il s'appuie sur le format Parquet, souvent utilisé pour stocker des
données de manière efficace, notamment dans des contextes d'analyse.

L'objectif était d'identifier les informations disponibles dans ce format et de
voir comment les exploiter pour produire les métadonnées attendues par Scan. Le
traitement devait notamment retrouver les informations géographiques, les champs
et les éléments nécessaires à la description de la ressource.

Pour ce format, le travail a porté sur la documentation. Il fallait comprendre
comment les informations spatiales étaient stockées dans le fichier et comment
retrouver les éléments utiles : champs, géométrie, système de coordonnées et
emprise. Ce format a nécessité une phase d'expérimentation, car il ne se
présente pas de la même manière qu'un fichier SIG traditionnel.

Les informations utiles ne sont pas toujours exposées avec les mêmes conventions
que dans d'autres formats. Le traitement devait donc rester prudent dans la
lecture des métadonnées et dans l'interprétation des informations géographiques.

La prise en charge de GeoParquet a donc permis d'élargir le périmètre de la
migration vers des formats plus modernes, tout en conservant la même exigence :
produire des résultats exploitables et comparables dans Scan.

% -----------------------------------------------------------------------------
% 3.2.4 Données CAO
% -----------------------------------------------------------------------------

\subsection{Données CAO}
\label{subsec:donnees_cao}

Une autre étape a concerné les données CAO. Ces données proviennent du dessin
assisté par ordinateur et ne suivent pas toujours les mêmes logiques que les
données SIG classiques.

Le travail sur ces formats a d'abord été engagé avec GDAL et OGR, comme pour les
premiers formats étudiés. Cette piste semblait cohérente, car ces bibliothèques
permettaient déjà de lire plusieurs types de données géographiques.

Cependant, cette approche a rapidement montré ses limites. Le format
\texttt{DWG}, important pour les données CAO, n'était pas pris en charge par
GDAL et OGR. Il a donc fallu rechercher une autre solution pour pouvoir exploiter
ces données dans le cadre de Scan.

Après cette phase d'étude, la solution retenue a été \textit{LibreDWG}. Elle a
permis de traiter les données CAO et d'en extraire les informations réellement
exploitables pour la documentation. Ces formats peuvent contenir des structures
particulières, des géométries différentes ou des informations moins directement
normalisées.

Certaines limitations dépendaient des technologies utilisées. Les formats
\texttt{SDF} et \texttt{DGN} ont notamment dû être mis de côté, car leur prise en
charge n'était pas suffisamment fiable dans le contexte de la migration.

Pour les données CAO, le travail est donc resté centré sur la documentation.
L'objectif était d'extraire les informations réellement disponibles et utiles à
Scan, sans chercher à forcer une interprétation SIG complète lorsque le format ne
s'y prêtait pas. Cette étape a montré qu'une migration ne consiste pas toujours à
tout reprendre immédiatement. Dans certains cas, il est préférable d'identifier
clairement les limites, de documenter les raisons du blocage et de concentrer les
efforts sur les formats pouvant être traités de manière fiable.

% =============================================================================
% 3.3 DEVELOPPEMENTS SUR LES BASES DE DONNEES
% =============================================================================

\section{Développements sur les bases de données}
\label{sec:bases_formats_avances}

Pour les bases de données, le choix a été différent de celui retenu pour les
formats fichiers. Nous avons privilégié l'utilisation de bibliothèques natives,
basées sur l'exécution de requêtes SQL. Cette approche permettait de s'appuyer
sur les capacités propres des bases de données et de bénéficier de leurs
performances, notamment lorsque les volumes à traiter étaient importants.

% -----------------------------------------------------------------------------
% 3.3.1 Oracle
% -----------------------------------------------------------------------------

\subsection{Oracle}
\label{subsec:oracle}

Après les premiers travaux sur les fichiers, le développement s'est poursuivi
sur Oracle. Cette technologie introduisait des contraintes différentes, car les
données sont stockées dans une base relationnelle, avec des paramètres de
connexion, des schémas, des tables et des droits d'accès.

Le traitement devait donc gérer la connexion à la base, identifier les tables
accessibles et récupérer les informations utiles à la documentation. Il fallait
aussi tenir compte des particularités d'Oracle Spatial pour les données
géographiques.

Pour Oracle, la bibliothèque \textit{python-oracledb} a été utilisée. Elle
permettait d'exécuter directement les requêtes nécessaires pour l'inventaire, la
documentation et la signature. Cette solution évitait de rapatrier inutilement
les données côté Python lorsque la base pouvait réaliser elle-même certaines
opérations.

Cette étape a demandé une attention particulière sur la robustesse. Une base de
données peut être inaccessible, une table peut ne pas contenir de géométrie ou
un utilisateur peut ne pas avoir les droits nécessaires. Le traitement devait
donc produire un comportement clair dans ces situations.

Le travail sur Oracle a renforcé l'idée que chaque format ou technologie possède
ses propres contraintes. Même lorsque les métadonnées recherchées sont proches,
la manière de les obtenir peut changer fortement selon la technologie.

% -----------------------------------------------------------------------------
% 3.3.2 SQL Server
% -----------------------------------------------------------------------------

\subsection{SQL Server}
\label{subsec:sql_server}

Le développement a également concerné Microsoft SQL Server. Comme pour Oracle,
il s'agit d'une base relationnelle pouvant contenir des données spatiales.

La prise en charge de cette technologie impliquait de gérer les paramètres de
connexion, les bases, les schémas et les tables accessibles. Il fallait aussi
tenir compte des types géographiques disponibles dans SQL Server et de la façon
dont les informations spatiales peuvent être interrogées.

Pour SQL Server, le choix s'est porté sur \textit{SQLAlchemy}. Cette bibliothèque
permettait de structurer les accès à la base tout en conservant une logique
reposant sur des requêtes SQL. Comme pour Oracle, l'objectif était de profiter
des capacités de la base pour limiter les traitements inutiles côté application.

Le travail a consisté à adapter les traitements aux spécificités de cet
environnement, tout en conservant une sortie cohérente avec les autres formats
et technologies.
L'objectif restait le même : permettre à Scan de recenser ou documenter les
ressources sans dépendre des anciens traitements.

Cette étape a aussi permis de renforcer la gestion des erreurs et des cas
particuliers. Les bases de données apportent souvent plus de variabilité que les
fichiers, car leur contenu dépend de l'organisation mise en place chez chaque
client.

% =============================================================================
% 3.4 TRAVAUX SUR LE NUAGE DES POINTS
% =============================================================================

\section{Traitement des nuages de points}
\label{sec:nuage_des_points}

Les travaux les plus récents portent sur le Nuage des points. Ce type de donnée
est particulier, car il peut représenter des volumes très importants et ne se
manipule pas comme une couche vectorielle classique.

L'objectif est d'identifier les métadonnées pertinentes pour décrire ces
ressources dans Scan. Il faut notamment comprendre les informations disponibles,
les formats utilisés et les limites liées au volume des données.

Cette partie est encore en cours. Elle nécessite de prendre en compte les
contraintes de performance, car les fichiers peuvent être très lourds. Le
traitement doit donc éviter les opérations trop coûteuses lorsque seules
certaines métadonnées sont nécessaires.

Le Nuage des points illustre bien la logique progressive de la migration. Chaque
nouveau format apporte ses propres contraintes et oblige à adapter les
traitements sans perdre la cohérence générale de la solution.

% Image à inclure :
% Une capture ou un schéma simple représentant une donnée de type Nuage des points.

% =============================================================================
% 3.5 PACKAGING ET EXECUTABLE AU FIL DES DEVELOPPEMENTS
% =============================================================================

\section{Packaging et exécutable au fil des développements}
\label{sec:packaging_fil_developpements}

% -----------------------------------------------------------------------------
% 3.5.1 Réutilisation de l'environnement Docker
% -----------------------------------------------------------------------------

\subsection{Réutilisation de l'environnement Docker}
\label{subsec:reutilisation_docker}

La génération de l'exécutable n'a pas été traitée comme une étape finale
indépendante. Une fois l'image Docker de construction mise en place, elle a été
réutilisée au fur et à mesure de l'avancement des scripts.

Cette organisation permettait de vérifier régulièrement que les développements
pouvaient toujours être intégrés dans un exécutable Windows. Elle évitait de
découvrir trop tard un problème lié au packaging ou à une dépendance.

Docker restait le seul prérequis pour générer l'exécutable. L'environnement de
construction contenait les versions nécessaires de Python, PyInstaller, GDAL et
des autres bibliothèques utilisées par Scan.

% -----------------------------------------------------------------------------
% 3.5.2 Vérification de l'exécutable
% -----------------------------------------------------------------------------

\subsection{Vérification de l'exécutable}
\label{subsec:verification_executable}

Après l'ajout ou la modification de certains traitements, l'exécutable pouvait
être régénéré afin de vérifier son bon fonctionnement. Les tests portaient alors
sur l'exécution du programme, la présence des dépendances nécessaires et la
cohérence des résultats produits.

Une attention particulière était portée aux bibliothèques géospatiales, car
elles reposent sur des dépendances natives. Le packaging devait donc embarquer
les éléments nécessaires pour éviter de dépendre d'une installation locale sur
la machine utilisée.

Cette vérification régulière permettait de garder un lien direct entre les
scripts développés et leur utilisation réelle dans Scan. Elle contribuait aussi
à sécuriser progressivement la migration vers une solution Python autonome.
